\documentclass[tikz]{standalone}
\usepackage{geometricfigures}

\usetikzlibrary{fit}

\usepackage{neuralnetwork}

\newcommand{\symplecticlayer}[1][] { \layer[bias=true,nodeclass={symplectic neuron},#1] }
\newcommand{\autoencoderlayer}[1][] { \layer[bias=true,nodeclass={autoencoder neuron},#1] }

\newcommand{\xin}[2]{$x_#2$}
\newcommand{\xout}[2]{$\hat x_#2$}

\begin{document}

\begin{neuralnetwork}[height=8, layertitleheight=30, toprow=true]

  % The package's styles, as this figure needs them: a neuron's text reads on its colour fill
  % in both themes, and a link title sits low with an almost transparent ground.
  \tikzset{
    neuron/.append style={text=onfill},
    linkstitle/.append style={text=gray, fill opacity=0.05, yshift=-1.8cm},
  }

  \tikzstyle{input neuron}=[neuron, fill=morange];
  \tikzstyle{hidden neuron}=[neuron, fill=mred];
  \tikzstyle{output neuron}=[neuron, fill=mgreen];
  \tikzstyle{symplectic neuron}=[neuron, fill=mblue];
  \tikzstyle{autoencoder neuron}=[neuron, fill=mpurple];

  \inputlayer[count=4, bias=false, text=\xin]

  \autoencoderlayer[count=6, bias=false]
  \linklayers[title=$\Psi^\mathrm{up}$]

  \hiddenlayer[count=6, bias=false]
  \linklayers%[title=$\psi^2_\mathrm{symp}$]
  % The transformer box encloses layer 1 (six neurons) and layer 2 (six neurons).
  \node[fit=(L1-1) (L1-6) (L2-1) (L2-6), draw, ultra thick, rounded corners, color=fg,
        label=above:Transformer] {};

  \outputlayer[count=4, text=\xout]%, title=Output Layer
  \linklayers[title=$\Psi^\mathrm{down}$]
  % \node[fit=\thislayer,draw, ultra thick, rounded corners, label=left:Transformer] (transformer) {};
\end{neuralnetwork}
\end{document}
